\section{Data and methods}
\label{sec:methodology}

The analysis proceeds in three stages. I first attach an occupation-level measure of exposure to artificial intelligence to each worker in a representative UK household survey. I then translate an assumed economy-wide AI shock into person-level changes to employment, wages and capital income, under alternative assumptions about who bears the displacement. Finally, I pass baseline and shocked data through an open-source tax-benefit microsimulation model to recover the distributional and fiscal consequences. This section describes each stage; the incidence scenarios that form the paper's central axis are formalised in Section~\ref{sec:incidence}.

\subsection{Measuring AI exposure}
\label{sec:exposure}

My starting point is the AI Occupational Exposure (AIOE) index of \citet{felten2021}, which links progress in ten AI application areas to the occupational ability requirements catalogued in O*NET, yielding a standardised score for each occupation that captures how much of its ability bundle overlaps with what AI systems can now do. Exposure alone does not distinguish occupations in which AI substitutes for workers from those in which it augments them. I therefore combine the AIOE with the complementarity adjustment of \citet{pizzinelli2023}, who construct an index $\theta_l \in [0,1]$ from O*NET work-context descriptors --- responsibility for others, physical proximity requirements, the consequence of error --- and job-zone preparation requirements. High-$\theta$ occupations are those in which AI is likely to complement human judgement rather than displace it. The complementarity-adjusted exposure measure (C-AIOE) shields high-complementarity occupations from the raw exposure score:
\begin{equation}
\label{eq:caioe}
\text{C-AIOE}_l = \text{AIOE}_l \times \bigl(1 - (\theta_l - \theta_{\min})\bigr),
\end{equation}
where $\theta_{\min}$ is the minimum complementarity across occupational groups, so that the least complementary group retains its full exposure score. Alternative exposure measures include \citet{tolan2021}, the task-level rubric of \citet{eloundou2023}, and the cross-country applications of \citet{cazzaniga2024}; Section~\ref{sec:results} reports the sensitivity of my headline results to substituting two of these for the C-AIOE.

Constructing the index for the UK requires a bespoke crosswalk, whose provenance I state in full because no published UK C-AIOE series exists at the level I need. AIOE scores are taken from the mapping to UK SOC2020 published with the DSIT analysis of AI and UK jobs \citep{dsit2023}, itself derived from the \citet{felten2021} scores via a chained crosswalk (US SOC2018 to US SOC2010, to ISCO-08, to UK SOC2020). The complementarity indices $\theta$ are not published at occupation level, so I reconstruct them from the public O*NET 27.3 release following the recipe documented by \citet{pizzinelli2023}; the reconstruction reproduces the published cross-occupation ordering up to a level offset, which is immaterial because equations~\eqref{eq:caioe} and~\eqref{eq:wage} use $\theta$ only relative to its minimum and its mean. Both series are aggregated to the nine SOC2020 major groups using employment weights from the Annual Survey of Hours and Earnings (ASHE) 2025 Table~14, which covers 340 of the 412 four-digit unit groups; the remainder carry no ASHE employment estimate and are excluded from the weighting. Finally, the C-AIOE scores are shifted so that the least-exposed major group scores exactly zero. Under the allocation rule in Section~\ref{sec:shocks} that group receives no job losses, and the shift prevents negative values of the standardised index from corrupting the allocation weights.

The one-digit level of aggregation is a genuine constraint. The Family Resources Survey (FRS) person records I use carry only the one-digit SOC2020 code, so exposure varies across just nine major groups and all within-group variation is lost. The exposure gradient I impose on the population is correspondingly coarse, and results should be read with that in mind; imputation of finer occupation codes from the Labour Force Survey is a planned refinement.

\subsection{Shock scenarios}
\label{sec:scenarios}

The size of the aggregate shock cannot be estimated from historical data and must be assumed; I treat it as a scenario input and vary it systematically. The central calibration takes a displacement rate of 7 per cent of employees and an aggregate wage uplift of 2.6 per cent for those who remain in work. Both figures are calibration conventions derived from \citet{briggs2023}: the 7 per cent converts their task-exposure estimates into an employment displacement share, and the 2.6 per cent converts their productivity projections into an aggregate wage gain. Neither is a forecast, and I do not present them as one. Alongside the labour-market effects, AI adoption is assumed to raise the return to capital by 0.4 percentage points, from a baseline of 1.005 per cent to 1.405 per cent.

Because displacement and wage growth are genuinely uncertain, the central calibration sits inside a grid spanning displacement rates from 0 to 10 per cent and aggregate wage growth from 0 to 5 per cent --- 66 scenario cells in all, with the capital-return shock applied throughout. The 1 per cent ``low'' anchor is a pure sensitivity case: \citet{acemoglu2025} estimates a ten-year GDP effect of around 1 per cent, not an employment displacement rate, and I do not attribute the anchor to him. Beyond the top of the grid, the 13 per cent ``high'' preset stress-tests a displacement rate roughly double the central convention.

\subsection{Implementing the shock at the micro level}
\label{sec:shocks}

\paragraph{Employment shock.} The aggregate number of displaced workers is the scenario displacement rate multiplied by weighted employee employment. The universe is all employees: workers whose SOC code cannot be matched to an exposure score form a pseudo-group assigned the employment-weighted mean exposure, so no employee is mechanically excluded from risk. The total is allocated across occupational groups $l$ in proportion to each group's employment multiplied by its exposure:
\begin{equation}
\label{eq:jobloss}
\text{JobLoss}_l = \text{TotalJobLoss} \times \frac{\text{EMP}_l \times \text{C-AIOE}_l}{\sum_{l} \text{EMP}_l \times \text{C-AIOE}_l},
\end{equation}
so job losses concentrate in large, highly exposed groups and the zero-scored least-exposed group loses none. Within each group, displaced individuals are selected by ordering group members on independent uniform random keys and filling the group's weighted quota from equation~\eqref{eq:jobloss}, with the record straddling the quota boundary included with probability equal to the remaining shortfall over its weight. Because the ordering keys are uniform, a record's inclusion probability does not depend on its grossing weight; the weights enter only through the quota accounting, which is the design intent.

Displacement is modelled as full withdrawal from employee work. A displaced worker's employment income, hours, employee pension contributions and salary-sacrifice amounts are set to zero, any statutory pay is removed, and labour-market status is set to unemployed before the shocked simulation runs. Dual jobholders retain their self-employment income; only the employee job is lost. PolicyEngine UK takes no unemployment-duration input, so displaced workers are assessed as current-period unemployed under the model's standard benefit rules rather than as long-term unemployed with contributory entitlements exhausted. To the extent that new-style (contributory) Jobseeker's Allowance would in practice be exhausted within a year, my results may modestly overstate the benefit support reaching displaced workers; I flag this in Section~\ref{sec:limitations}.

\paragraph{Wage shock.} Each worker who survives the employment shock receives a percentage uplift proportional to complementarity, normalised by the employment-weighted mean complementarity over baseline workers:
\begin{equation}
\label{eq:wage}
\text{WageChange}_l = \text{WageGrowth} \times \frac{\theta_l}{\bar{\theta}} \times \text{Wage}_l,
\qquad
\bar{\theta} = \frac{\sum_l \text{EMP}_l \, \theta_l}{\sum_l \text{EMP}_l},
\end{equation}
where $\text{EMP}_l$ is baseline (pre-shock) employment. This is the baseline wage rule. Two properties should be noted. The normalisation makes the \emph{employment-weighted mean} percentage uplift equal the assumed wage growth. The aggregate wage bill grows by exactly that rate only if $\theta$ is uncorrelated with pay. Since high-$\theta$ occupations earn more, the wage-bill-weighted uplift is somewhat larger. A second property: because $\bar{\theta}$ is computed over all baseline workers while the uplift accrues only to survivors---who, being drawn from less-exposed occupations, have above-average $\theta$---displacement raises the realised mean survivor uplift slightly above the headline rate. Exact conservation is verified in seeded tests at zero displacement; both departures are properties of the rule itself rather than implementation error. The economic logic mirrors equation~\eqref{eq:caioe}: AI complements rather than substitutes for high-$\theta$ occupations, so those occupations capture a larger share of the productivity gains --- an assumption the compression scenario of Section~\ref{sec:incidence} deliberately inverts.

\paragraph{Capital shock.} The 0.4 percentage-point increase in the return to capital is implemented by scaling each person's savings-interest and dividend income by the ratio of shocked to baseline return, $1.405/1.005 \approx 1.398$. The baseline return of 1.005 per cent is taken from the baseline calibration; because the scaling factor is the ratio of two return levels, the size of the capital channel is highly sensitive to this assumption. The same $+0.4$ percentage-point shock applied to, say, a 4 per cent baseline return would scale capital incomes by only 10 per cent rather than 40 per cent. The capital-channel results should therefore be read as conditional on this calibration; equivalently, the shock can be read directly as a 39.8 per cent proportional increase in interest and dividend income. Rental income is excluded from the scaling, on the modelling assumption that the AI productivity channel operates through financial rather than housing capital. The self-employed are outside both labour-market shocks; their interest and dividend income is scaled like everyone else's.

\subsection{Incidence scenarios}
\label{sec:incidence}

The rules above fix how large the shock is; they do not settle who bears it, and the empirical literature has not settled that question either. I therefore treat incidence as a first-class scenario axis. Holding the central aggregate shock fixed --- 7 per cent displacement, a 2.6 per cent wage uplift, the capital shock --- I re-allocate the same weighted quota of displacement under four incidence families:

\paragraph{Exposure-proportional.}The default allocation of equations~\eqref{eq:jobloss} and~\eqref{eq:wage}: displacement follows the C-AIOE gradient and wage gains follow complementarity. This family is the central calibration used throughout the shock-size grid.

\paragraph{Junior-concentrated.}Displacement draws are tilted towards workers aged 16--24 by a multiplier equal to the ratio of junior to overall displacement effects reported by \citet{kleinteeselink2025} for the UK ($5.8/4.5$), consistent with the seniority-biased evidence of \citet{hosseini2026} and \citet{brynjolfsson2025}. The group quotas and wage rule are otherwise unchanged.

\paragraph{Expertise-compression.}A stylised, author-designed stress test of the view that AI compresses expert rents rather than displacing routine work. Displacement draws are tilted (multiplier 2) towards the top weighted-earnings tertile of workers in above-median-exposure occupations, and the wage uplift is allocated by \emph{inverse} complementarity --- $\theta$ is reflected about its range, $\theta'_l = \theta_{\max} + \theta_{\min} - \theta_l$, so that mid-skill workers capture the capability gains --- with the same employment-weighted normalisation as equation~\eqref{eq:wage}.

\paragraph{Uniform.}Every employee is assigned equal exposure and equal complementarity, so displacement risk and wage gains are flat across occupations. This family isolates the contribution of the incidence gradient itself: differences between it and the others are attributable entirely to who bears the shock.

Each family delivers the same weighted count of displaced workers to within rounding, so cross-family differences in fiscal cost, poverty and inequality measure the consequences of incidence alone. The incidence gradient in each family is an assumption, not an estimate; the paper's contribution is to show how the tax-benefit system mediates whichever incidence materialises.

\paragraph{Measured-incidence family.}Section~\ref{sec:results-incidence} adds a fifth family that replaces the stylised gradients above with the displacement patterns measured in UK firm-level data by \citet{kleinteeselink2025}. Displacement probability is proportional to the shifted C-AIOE multiplied by three tilts: a junior multiplier of 1.29 (the ratio of their junior to overall displacement effects, $5.8/4.5$), a London multiplier of 1.5, and a wage-tier multiplier of 9.6 above versus 0.6 below weighted-median earnings, reflecting their finding that losses concentrate in high-paying firms; the aggregate shock stays at the central preset. The reader should treat this family as an author-constructed, top-loaded stress test loosely anchored to \citet{kleinteeselink2025} rather than an implementation of it.\footnote{Verification against the full text (SSRN working paper, December 2025 revision) demands more candour than a citation conveys. The paper's own estimates are difference-in-differences semi-elasticities per standard deviation of LLM exposure: total employment $-0.3$ per cent, junior employment $-0.4$ per cent (a ratio close to the 1.29 used here, though built from two rounded coefficients of which the senior effect is statistically insignificant), and high-compensation firms $-0.7$ per cent against a statistical zero for low-compensation firms. The 9.6/0.6 wage-tier figures descend from the accompanying press release's scaling to maximally exposed firms, not from the paper's tables, and the pay split is by \emph{firm average compensation}, which I proxy by person earnings. Most seriously, the London tilt of 1.5 is \emph{contradicted} by the revised full text, in which London-headquartered firms show minimal, statistically insignificant employment effects while the losses concentrate outside London; I retain the tilt so that the family matches its originally circulated form. Section~\ref{sec:results-incidence} shows the family's fiscal and poverty results are robust to replacing 9.6/0.6 with much milder ratios; the multiplicative independence of the three channels and the age-under-25 proxy for their seniority concept remain author-imposed throughout.}

\subsection{Microsimulation with PolicyEngine UK}
\label{sec:microsim}

I use PolicyEngine UK (version 2.89.2), an open-source tax-benefit microsimulation model, running on PolicyEngine UK's processed (base, unenhanced) Family Resources Survey 2024--25 microdataset. This reshapes the UK Data Service FRS 2024--25 (study SN~9563) into PolicyEngine's tax-benefit schema, applies source-level imputations (council tax, benefit take-up, salary sacrifice) and retains the survey's own grossing weights; I obtain it from PolicyEngine's \texttt{policyengine-uk-data} repository and load it through \texttt{UKSingleYearDataset} and \texttt{Microsimulation}. Occupation codes are not carried in the processed dataset, so I join the SOC2020 major group from the raw FRS \texttt{adult.tab} file (UK Data Service SN~9563) onto the simulation's person table using the identifier convention $\text{person\_id} = \text{SERNUM} \times 1000 + \text{PERSON}$, which matches every FRS adult. The simulation year is 2026.

The counterfactual design is a paired comparison on identical data: a baseline simulation and a shocked simulation are run on the same records, with the shocked run receiving the modified employment income, hours, pension contributions, savings-interest income, dividend income and employment status of Section~\ref{sec:shocks} as direct inputs. All reported effects are differences between the two runs. This design removes between-run noise --- the survey sample and grossing weights are identical on both sides, so no part of a reported effect reflects resampling --- but it does not quantify population sampling uncertainty: both runs inherit whatever sampling error the FRS itself carries, and my Monte Carlo intervals speak only to the stochastic identity of the displaced, not to survey sampling.

All income results are expressed on the HBAI cash disposable income concept --- equivalised household net income as defined for the Households Below Average Income statistics and implemented in PolicyEngine UK --- which is also the concept underlying the poverty lines. I report four outcome families: the change in the government balance (the Exchequer cost of the shock, net of additional revenue on higher wages and capital income); person-weighted poverty rates before and after housing costs (BHC and AHC); the Gini coefficient of equivalised household disposable income, computed with negative incomes bottom-coded at zero; and mean changes in household income by decile, with deciles fixed at their baseline positions so that individuals are tracked rather than re-ranked.

Because the identity of the displaced is stochastic, I adopt explicit seed conventions and label every result accordingly. The 66-cell grid, the three presets and the decomposition figure are single draws at seed~0, so that cells differ only in their scenario inputs, never in their random draw. The five incidence families report means over 20 draws (seeds 0--19). The decile transition and wage profiles (Figures~\ref{fig:transition} and~\ref{fig:wages}) average 50 seeded draws. The robustness summary and the decomposition confidence band use a 20-draw Monte Carlo at seeds 0--19; wherever a $\pm$ value appears it is the standard deviation across draws, not a survey-based standard error.

\paragraph{Policy reforms in the shocked world.}The reforms of Section~\ref{sec:policy} are evaluated as $M(\text{shock}+\text{reform}) - M(\text{shock})$ on the fixed seed-0 shock draw, applying for 2026 only, with poverty lines held at their shocked-world values. Benefit-side reforms are implemented as PolicyEngine UK parameter reforms --- scaling the Universal Credit (UC) standard allowance, raising the benefit-cap thresholds and cutting the taper rate --- so their costs are net of the means-testing clawback the model computes. Wage insurance is instead a direct post-simulation transfer to displaced workers (half of lost earnings, capped), treated as non-taxable and disregarded by means tests, so its gross cost equals its net cost and brackets the generous end of possible implementations.

\paragraph{Constituency geography.}The constituency results of Section~\ref{sec:results-geo} rest on PolicyEngine's $650 \times 53{,}508$ weight matrix, which assigns each household in the \emph{enhanced} FRS 2023--24 dataset a calibrated weight in every Westminster parliamentary constituency. The SERNUM-based occupation join of Section~\ref{sec:microsim} is invalid on that dataset, so SOC major group is imputed from the plain-FRS 2024--25 weighted distribution within age band $\times$ gender $\times$ region $\times$ earnings-decile cells (99.4 per cent weighted coverage from the full four-way cell, seed 0). Constituency variation therefore reflects demographic and earnings composition plus a single displacement draw, not observed occupation mix, and the geographic results --- computed on a different dataset and period --- are not directly comparable to the poverty numbers elsewhere in the paper.
